% Part: many-valued-logic
% Chapter: syntax-and-semantics
% Section: semantic-notions

\documentclass[../../../include/open-logic-section]{subfiles}

\begin{document}

\olfileid[pt]{mvl}{syn}{sem}

\olsection{Noções semânticas}

Suponha que uma lógica multivalorada $\Log L$ seja dada por uma matriz. Então
podemos definir as noções semânticas usuais para~$\Log L$.

\begin{defn} 
\begin{enumerate}
\item !!^a{formula}~$!A$ é \emph{satisfatível} se, para
  algum~$\pAssign{v}$, $\pSat{v}{!A}$; ela é
  \emph{insatisfatível} se não há $\pAssign{v}$ tal que $\pSat{v}{!A}$;
\item !!^a{formula}~$!A$ é uma \emph{tautologia} se $\pSat{v}{!A}$ para
  todas as !!{valuation}s~$v$;
\item Se $\Gamma$ é um conjunto de !!{formula}s, $\Gamma \Entails !A$ (``$\Gamma$
  tem $!A$ como consequência'') se, e somente se, $\pSat{v}{!A}$ para toda
  !!{valuation}~$\pAssign{v}$ para a qual $\pSat{v}{\Gamma}$.
\item Se $\Gamma$ é um conjunto de !!{formula}s, $\Gamma$ é
  \emph{satisfatível} se existe !!a{valuation}~$\pAssign{v}$ para a qual
  $\pSat{v}{\Gamma}$, e $\Gamma$ é
  \emph{insatisfatível} caso contrário.
\end{enumerate} 
\end{defn}

Para essas noções, valem alguns dos mesmos fatos que no
caso da lógica clássica:

\begin{prop}
\ollabel{prop:semanticalfacts} 
\begin{enumerate} 
\item $!A$ é uma tautologia se, e somente se,
  $\emptyset \Entails !A$; 
\item Se $\Gamma$ é satisfatível, então todo subconjunto finito de $\Gamma$
  também é satisfatível; 
\item\ollabel{def:monotonicity}%
Monotonicidade: se $\Gamma \subseteq \Delta$
  e $\Gamma \Entails !A$, então também $\Delta \Entails !A$;
\item\ollabel{def:Cut}%
Transitividade: se $\Gamma \Entails !A$ e
  $\Delta \cup \{ !A\} \Entails !B$, então $\Gamma \cup \Delta \Entails
  !B$;
\end{enumerate}
\end{prop}

\begin{proof}
Exercício.
\end{proof}

\begin{prob}
Demonstre \olref[mvl][syn][sem]{prop:semanticalfacts}.
\end{prob}

Na lógica clássica, podemos relacionar a consequência e o condicional. Por
exemplo, temos a validade do \emph{modus ponens}: se $\Gamma
\Entails !A$ e $\Gamma \Entails !A \lif !B$, então $\Gamma \Entails
!B$. Outra relação importante entre $\Entails$ e $\lif$ na
lógica clássica é o teorema semântico da dedução: $\Gamma \Entails !A
\lif !B$ se, e somente se, $\Gamma \cup \{!A\} \Entails !B$. Esses
resultados \emph{nem sempre} valem em lógicas multivaloradas. Sua validade
depende da função de verdade~$\tf{\lif}$.

\end{document}
